\chapter{Formalization Roadmap}

\section{What counts as progress}

The project should prefer theorem statements that end in one of the following
finite conclusions:
\[
  \RealRaw.\Valid,\qquad
  \ComplexRaw.\Valid,\qquad
  x\Equiv y,\qquad
  \close_\eps(I,J).
\]
These are the computational substitutes for classical existence and equality.
Whenever a classical proof says "choose a sufficiently fine partition" or
"take a small enough neighborhood", the blueprint should identify the
corresponding explicit modulus.

\section{Near-term milestones}

\begin{enumerate}
  \item Prove the exact rational FTC route:
  turn \cref{thm:finite-ftc-estimate} into Lean constructors for
  `EffectiveFTCExact` and `EffectiveFTC`.

  \item Prove integrability from interval regularity:
  refine \cref{conj:integral-existence} into concrete Riemann-sum plans,
  nesting lemmas, and width bounds.

  \item Extend polynomial calculus:
  formalize exact derivative and FTC certificates for all rational
  polynomials, then rational functions on certified denominator-nonzero
  intervals.

  \item Finish arctangent soundness:
  prove that geometric arctangent and power-series arctangent agree on
  \([-1,1]\), giving \cref{conj:leibniz-machin}.

  \item Finish geometric pi validity:
  prove dyadic nesting and width shrinkage for the area and circumference
  algorithms, completing \cref{conj:geometric-pi-validity}.

  \item Choose an FTA route:
  either develop algebraic complex numbers with isolation boxes, or develop
  the argument-principle box-search route.
\end{enumerate}

\section{Lean declaration index}

\begin{itemize}
  \item Foundations:
  \code{QInterval}, \code{QBox}, \code{RealRaw}, \code{RealRaw.Equiv}, \code{ComplexRaw},
  \code{RealFunRaw}, \code{FunctionOnInterval}, \code{IntervalRegularOn}.

  \item FTC:
  \code{EffectiveDerivativeExact}, \code{EffectiveFTC},
  \code{FTC.effectiveFTC\_equiv\_endpoint},
  \code{DefiniteIntegralEqualsEndpointDifference}.

  \item FTA:
  \code{CPoly.Coeffs}, \code{IsComputableRoot}, \code{AlgebraicComplex},
  \code{ComputableFTA}, \code{AlgebraicFTA},
  \code{ComputableFTA\_of\_AlgebraicFTA}.

  \item Pi:
  \code{piLeibniz}, \code{piMachin}, \code{piCircleArea},
  \code{piCircumference}, \code{PiProofs.archimedesTheorem},
  \code{PiProofs.LeibnizEqMachin}.

  \item Series and elementary functions:
  \code{DirichletSeries.zetaTwoRaw},
  \code{Basel.eulerBasel\_geometricPi},
  \code{ExpProofs.EPowerSeriesEqEuler},
  \code{Taylor.IteratedFTCChain},
  \code{Taylor.ArctanKernel.finiteRemainderRoute}.
\end{itemize}

\section{A useful style rule}

Every major theorem should have three layers:
\begin{enumerate}
  \item a readable mathematical statement in the blueprint;
  \item a Lean proposition or structure with the same quantifier shape;
  \item small finite lemmas that discharge the numerical or algebraic
  estimates.
\end{enumerate}
This keeps the project faithful to the historical mathematics while making
the formal work genuinely computable.
